\clearpage
\begingroup
\let\clearpage\relax
\let\cleardoublepage\relax
\vspace*{\fill}
\chapter{软件工程概述与过程模型}
\begin{center}
软件工程用「工程化」的方法控制软件的成本、进度与质量。本章建立「软件危机--过程模型--生命周期--敏捷」四块基石。
\end{center}
\vspace*{\fill}
\endgroup
\clearpage

\section{软件与软件危机}

\begin{itemize}
    \item \textbf{软件}：程序 + 数据 + 文档。它是逻辑产品，无物理磨损，但会因变更而「退化」。
    \item \textbf{软件危机}：需求不清、进度失控、成本超支、质量低、维护难。
    \item \textbf{软件工程}：将系统化、规范化、可度量的工程方法应用于软件的开发、运行与维护。
\end{itemize}

\section{软件过程模型}

\begin{table}[H]
\centering
\caption{常见过程模型对比}
\begin{tabulary}{\textwidth}{CCC}
\toprule
模型 & 特点 & 适用 \\
\midrule
瀑布 (Waterfall) & 阶段线性、文档驱动 & 需求稳定、合同型 \\
迭代/增量 & 分轮交付，逐步完善 & 中大型、需求渐进 \\
螺旋 (Spiral) & 迭代 + 风险分析 & 高风险大型项目 \\
原型 (Prototype) & 先做可运行原型澄清需求 & 需求模糊 \\
敏捷 (Agile) & 短迭代、响应变化、可工作软件 & 需求多变、小团队 \\
\bottomrule
\end{tabulary}
\end{table}

\section{软件生命周期}

\begin{figure}[H]
\centering
\begin{tikzpicture}[font=\small, node distance=1.3cm]
    \node[draw, rounded corners] (req) {需求};
    \node[draw, rounded corners, right=of req] (design) {设计};
    \node[draw, rounded corners, right=of design] (impl) {实现};
    \node[draw, rounded corners, right=of impl] (test) {测试};
    \node[draw, rounded corners, right=of test] (maint) {维护};
    \draw[->] (req) -- (design);
    \draw[->] (design) -- (impl);
    \draw[->] (impl) -- (test);
    \draw[->] (test) -- (maint);
\end{tikzpicture}
\caption{软件生命周期的主要阶段}
\end{figure}

\noindent 维护成本常占生命周期总成本的 $60\%\sim80\%$；需求与设计缺陷的修复代价随阶段指数增长：
\begin{equation}
    C_{\text{fix}}(\text{阶段}) \propto e^{\,k\cdot \text{阶段序号}}
\end{equation}

\section{敏捷与 Scrum}

\begin{itemize}
    \item \textbf{敏捷宣言}：个体与互动、可工作软件、客户合作、响应变化 高于 流程与工具、详尽文档、合同谈判、遵循计划。
    \item \textbf{Scrum}：Sprint（1--4 周）、产品待办列表 (Product Backlog)、冲刺待办列表、每日站会、评审与回顾。
    \item \textbf{XP}：结对编程、测试驱动开发 (TDD)、持续集成、小步发布。
\end{itemize}

\begin{table}[H]
\centering
\caption{Scrum 三种角色与三类工件}
\begin{tabulary}{\textwidth}{CC}
\toprule
角色 & 工件 \\
\midrule
产品负责人 (PO) & 产品待办列表 \\
Scrum Master & 冲刺待办列表 \\
开发团队 & 增量 (Increment) \\
\bottomrule
\end{tabulary}
\end{table}
